
\chapter{Introducción} 

La Universidad Nacional de Costa Rica (UNA) busca modernizar y optimizar sus sistemas de gestión, especialmente aquellos relacionados con la calidad académica y la administración de recursos. En este contexto, la universidad enfrenta desafíos significativos en dos áreas clave: por un lado, la gestión de evidencias para los procesos de acreditación ante el Sistema Nacional de Acreditación de la Educación Superior (SINAES), y por otro, la administración eficiente de los tiempos de jornada universitaria, un recurso vital para la docencia y los proyectos institucionales.

Actualmente, la gestión de evidencias SINAES se realiza de manera manual y se apoya en una estructura de almacenamiento basada en carpetas en Google Drive. Este enfoque genera ineficiencias significativas, como la redundancia documental, una clasificación manual que consume un promedio de dos horas diarias por coordinador, y una marcada dificultad para correlacionar las evidencias con los criterios específicos de calidad definidos por SINAES. El problema principal radica en la ausencia de un sistema automatizado que permita una categorización eficiente y la falta de metadatos asociados a los documentos, lo que resulta en una dependencia de criterios individuales para la clasificación. Estas deficiencias no solo conllevan un desperdicio de recursos institucionales y la duplicación de la memoria utilizada, sino que también incrementan el riesgo de errores en la verificación del cumplimiento de criterios y pueden generar retrasos críticos en la preparación de informes para auditorías SINAES. Además, la falta de seguridad adecuada para los documentos institucionales representa un riesgo potencial de pérdida de acreditación en carreras clave. Para mitigar estos desafíos, se propone el desarrollo de un módulo de gestión de evidencias como parte del macro-proyecto "Gestión de Calidad", el cual estará enfocado exclusivamente en la interfaz de usuario para la manipulación directa de documentos en Google Drive, sin una capa de backend dedicada en esta fase inicial.
